\scanpage{112}
\mspage{111}
\uline{Atelier}\inserted{(\unsure{spéciaux|speciaux})} : est donné par $(\unsure{\mathcal{F}|\mathscr{F}},\leq,\ll;\unsure{|\circ|\,|\,|o|})$ où $\mathcal{F}$ est un ensemble muni de deux relations d'ordre $\leq$, $\ll$, et d'une \deleted{\unsure{donnée|relation} \unsure{supplémentaire}} relation. $\deleted{\unsure{\gtrless|\lessgtr}}$ $\deleted{\unsure{|\circ|}}$ Axiomes

\uline{At 1} $\mathcal{F}$\inserted{$\ne\varnothing$ et $\mathcal{F}$} stable par $\mathrm{Sup}^{\leq}$ \deleted{de \unsure{f.}} \uline{majorés} quelconques.

\uline{At 2} Toute figure\inserted{$F$} est le $\mathrm{Sup}^{\leq}$ de \inserted{l'ens $\widetilde{F}$ de} ses ``strates''
\[
\widetilde{F}=\{X\in\mathcal{F}\mid X\leq F,\ X\text{ ``strictement irréductible''}\}
\]
% slanted
\margin{\unsure{équivalent à} : $F\leq G\Longleftrightarrow\widetilde{F}\subset\widetilde{G}$ (critère \unsure{des deux} figures par strates) $\widetilde{F}$, \deleted{\unsure{s'app|s'ap}} muni de $\leq$, s'appelle ``la \uuline{\unsure{maquette}} de la figure'' (qu'elle \unsure{restitue|reconstitue} par ``déploiement'').}

\deleted{\unsure{At 3}} Si $M$ désigne l'ens des figures élémentaires -- multistrates (objets str. irréd. de $\mathcal{F}$), $\deleted{\unsure{M}}$ muni de la structure $\leq$ seulement, \unsure{donc|alors} la donnée de $(\mathcal{F},\leq)$ satisfaisant \unsure{:|à} At 1, At 2 \deleted{\unsure{pure}} est équivalente à celle d'un ensemble ordonné $(M,\leq)$, avec un ens.\inserted{$\widetilde{\mathcal{F}}$} de parties fermées\inserted{de $M$} (les $\widetilde{F}$), contenant avec chaque $\unsure{\Phi|\Omega}$ ses parties fermées, \deleted{\unsure{et}} recouvrant $M$ \deleted{\unsure{les}\ $\unsure{\widetilde{F}}$} et contenant au moins la partie vide de $M$ (``figure vide''). Ceci posé
% slanted
\margin{Les $\mathrm{Sup}$\inserted{, $\mathrm{Inf}$} dans $\mathcal{F}$ s'expriment par $\cup$, $\cap$ dans $\widetilde{\mathcal{F}}$.}

% The remaining lower block is cancelled by long diagonal strokes.
\deleted{\uline{At 3}}
\[
\deleted{F\mathrel{\unsure{\ll|\leq}}G\Longleftrightarrow\forall X\in\widetilde{F},\ \exists Y\in\widetilde{G}\text{ avec }X\ll Y}
\]
\deleted{\inserted{et si $\forall X\in\widetilde{F}$, $Y\in\widetilde{G}$,}}
\deleted{\uline{$X\cup Y$ existe} dans $\mathcal{F}$ (i.e. $\{X,Y\}$ majoré pour $\leq$), alors $F\cup G$ \unsure{.|:} existe dans $\mathcal{F}$.}

\deleted{On introduit relation dans $\deleted{\unsure{M}}$ $\mathcal{F}$}
\[
\deleted{\unsure{F}\mathrel{\gtrless}\unsure{G}\Longleftrightarrow\{\unsure{F},\unsure{G}\}\text{ majoré dans }\mathcal{F}}
\]
\deleted{\inserted{i.e. $F\cup G$ existe}}
\deleted{Par \uline{At 3}, il suffit de la connaître dans $M$.}
